\documentclass[10pt,twocolumn]{article}

\usepackage[margin=0.67in,columnsep=0.22in]{geometry}
\usepackage[T1]{fontenc}
\usepackage{lmodern}
\usepackage{microtype}
\usepackage{amsmath,amssymb,mathtools}
\usepackage{booktabs,tabularx,array,multirow}
\usepackage{enumitem}
\usepackage{xcolor}
\usepackage{hyperref}
\usepackage{url}

\definecolor{linkblue}{RGB}{28,78,121}
\hypersetup{colorlinks=true,allcolors=linkblue,pdfauthor={Raina-ya},
  pdftitle={Runtime Visibility Limits of Bias Probes for Quantum Error-Mitigated Selection}}
\setlist{nosep,leftmargin=*}
\setlength{\parindent}{0.8em}
\setlength{\parskip}{0pt}
\setlength{\tabcolsep}{3.5pt}
\renewcommand{\arraystretch}{1.08}
\newcommand{\E}{\mathbb{E}}
\newcommand{\Prb}{\mathbb{P}}
\newcommand{\ind}{\mathbb{1}}
\newcommand{\code}[1]{\texttt{#1}}

\title{\textbf{Runtime Visibility Limits of Bias Probes for\\
Quantum Error-Mitigated Selection}\vspace{0.35em}\\
\large A Preregistered Fixed-Budget Study}
\author{Raina-ya\\Affiliation: TODO\\[0.3em]
\small Pre-results manuscript draft --- Protocol v4}
\date{13 November 2026}

\begin{document}
\maketitle

\begin{abstract}
Confidence intervals quantify sampling uncertainty but do not, by themselves,
bound differential systematic bias. A quantum decision can therefore be
reported with high sampling confidence while selecting the wrong candidate.
We formulate runtime bias diagnosis as a visibility problem: which probes can
detect a bias large enough to reverse a QAOA--MaxCut top-1 decision, which are
blind by construction, and when probe shots are worth removing from the target
measurement budget. Protocol v4 preregisters six bias probes against seven
non-null gate and readout structures (42 cells), with a fixed total budget of
16,384 shots per graph--cell--repeat. The confirmatory design uses 100 holdout
graphs and 32 repeats; analytic visibility strata, matched nulls, operating
rules, seeds, and decision endpoints were frozen before noisy development or
holdout evaluation. The primary endpoint is the change in false-certification
rate between a CI-only rule and a bias-aware abstaining rule at equal total
cost. This pre-results manuscript reports only analytic results, split-99
software-design checks, a documented defect in the superseded v3 protocol, and
a non-circular power calculation. No development or confirmatory holdout
result is reported in this manuscript snapshot, and no hardware or
physical-probe performance result is reported. The central contribution is
thus a falsifiable and auditable test of whether runtime diagnostics add
decision value rather than merely respond to noise.
\end{abstract}

\section{Introduction}

Quantum error mitigation (QEM) trades additional circuits, shots, or model
assumptions for reduced estimator error, without implementing full fault
tolerance. Canonical examples include zero-noise extrapolation and related
near-term mitigation methods\cite{temme2017,li2017,endo2018}. Their performance
is commonly summarized by estimator bias, variance, or mean-squared error.
Selection tasks impose a different requirement: only the sign of a contrast
between candidates determines the action. Improved average accuracy need not
imply a more reliable decision, particularly when the two candidates acquire
different systematic biases.

Let the ideal candidate values be $\mu_A$ and $\mu_B$, with ideal gap
$g=\mu_A-\mu_B$ and margin $\Delta=|g|$. Under an error process, let the
population gap be $g_{\mathrm n}=g+b_A-b_B$. The decision-relevant bias is
therefore
\begin{equation}
  \beta=(b_A-b_B)\operatorname{sign}(g).
  \label{eq:beta}
\end{equation}
The ideal decision is reversed whenever $\Delta+\beta<0$. A narrow interval
around $g_{\mathrm n}$ can exclude zero and still certify the wrong sign,
because the interval controls finite-shot variation around a biased center.

This study asks whether runtime probes can reveal such harmful differential
bias under a fixed shot budget. The question is deliberately narrower than
general device characterization. A useful probe must satisfy three conditions:
it must respond to the harmful direction, retain adequate finite-shot power,
and intercept enough wrong certifications to compensate for the target shots
it consumes. A probe that fails the first condition is blind in principle; a
probe that fails the second is visible only at the population level; and a
probe that fails the third can be diagnostically informative but harmful to
the deployed decision rule.

The framework draws on partial identification, where observations may define a
set rather than a point\cite{imbens2004}; negative controls, whose validity
requires a shared bias mechanism\cite{lipsitch2010}; and known learnability
limits for quantum noise\cite{chen2023learnability}. It also connects to
established gate, readout, and randomized-compiling tools
\cite{bravyi2021,maciejewski2021,wallman2016,czarnik2021}. The novelty claim is
not that any individual probe or concentration bound is new. It is the
preregistered joint test of analytic visibility, finite budget, and
decision-level false certification across a common QAOA selection task.

The study is simulation-based. No present claim establishes hardware utility,
and the synthetic three-body readout stress structure has no asserted hardware
provenance. Hardware experiments are a later go/no-go stage, conditional on
the simulation results and independently frozen procedures.

\section{Methodology}

\subsection{Decision, identification, and visibility}

For QAOA--MaxCut\cite{farhi2014}, a bit string $z$ on graph $G=(V,E)$ receives
the normalized score
\begin{equation}
  s_G(z)=\frac{1}{|E|}\sum_{(i,j)\in E}\ind\{z_i\neq z_j\}\in[0,1].
\end{equation}
Each graph has two frozen candidate parameter vectors. Their expected scores
define $g$; the sign of an observed sample gap defines the selected candidate.

Let $D$ denote all target and probe observations and $\mathcal M(D)$ the
admissible error mechanisms compatible with them. A generic identification
region for the ideal gap is
\begin{equation}
  \mathcal I_g(D)=\{g(M):M\in\mathcal M(D)\}.
  \label{eq:idregion}
\end{equation}
If the region contains both positive and non-positive values, the decision is
not sign-identified. This is a population concept, distinct from a confidence
interval for shot noise. The project uses it to state blindness, not to claim
that all 42 cells require explicit optimization of Eq.~\eqref{eq:idregion}.

For each probe--structure cell, Protocol v4 freezes a component statistic and
a matched null. A cell is \emph{invariant} if the population statistic is
identical under the structure and its matched null; no threshold can then
recover mechanism-specific power. A \emph{sensitive} cell has a nonzero
population contrast, but may remain underpowered at the frozen budget. The
visible/partial/blind label combines this analytic relation with the frozen
power and false-alarm rules. Invariant and sensitive cells are reported
separately so that construction-driven agreements cannot mask finite-budget
failures.

\subsection{Probes}

Table~\ref{tab:probes} summarizes the six primary probes. P0, the target-only
confidence-interval comparator, is intentionally excluded from the visibility
matrix because an unresolved target interval is not a bias alarm. Two further
statistics are recorded but are exploratory and carry no prediction: P1b, the
negative quasi-probability weight of the scale-1/3 Richardson distribution
built from P1 counts, and GVS, a covariance-survival ratio computed from target
counts. Both reuse existing counts at no extra shot cost. Positive-cost
probes receive $n_P=1{,}638$ shots and the targets receive $n_T=14{,}746$;
zero-extra-shot P5 retains all $N=16{,}384$ target shots.

\begin{table*}[t]
\centering
\caption{Each probe targets one projection of the error process, and the
preregistered matrix predicts it to be blind to most structures; no probe is
treated as an omnibus detector. Positive-cost probes take $n_P=1{,}638$ of the
$N=16{,}384$ shots; P5 reuses target counts. The last column lists the
structures whose cells are predicted blind (Table~\ref{tab:matrix}).}
\label{tab:probes}
\begin{tabularx}{\textwidth}{@{}lXlXl@{}}
\toprule
Probe & Frozen statistic & Extra cost & Physical target & Predicted blind under \\
\midrule
P1 & $|(g_3-g_1)/2|$ from local unitary folding at scales 1 and 3 & 2 circuits, $n_P$ & Noise that scales with gate count & E2, E4, E5, E5$'$, E6 \\
P2 & $|g_{\rm twin,obs}-g_{\rm twin,ideal}|$ for a deterministic nearest-Clifford twin & 2 circuits, $n_P$ & Gate bias transferable to the twin & E5$'$ \\
P3 & $|g_{\rm plain}-\overline{g}_{\rm frame}|$ over 16 Pauli-frame versions per candidate & 32 circuits, $n_P$ & Coherent error around two-qubit gates & E1, E2, E4, E5, E5$'$, E6 \\
P4a & Single-qubit assignment deviation from all-zero/all-one preparations; temporal contrast for E6 & 2 circuits, $n_P$ & Readout marginals and drift & E1, E2, E3, E5, E5$'$ \\
P4b & One frozen component per cell: marginal, pair-correlation, or temporal contrast from pair preparations & $4\binom{w}{2}$ circuits, $n_P$ & Pairwise readout correlation & E1, E2, E3, E5$'$ \\
P5 & Bit-complement symmetry distance of target counts & none & Flip-symmetry-breaking error & E1, E2, E5, E5$'$ \\
\bottomrule
\end{tabularx}
\end{table*}

Within each width--depth stratum and cell component, development matched-null
scores $T_{(1)}\le\cdots\le T_{(n)}$ determine
\begin{equation}
  \tau=T_{\min\{n,\lceil(n+1)(1-\alpha)\rceil\}},\qquad \alpha=0.05,
\end{equation}
and the alarm is strict: $T>\tau$. Decision endpoints instead calibrate alarms
against clean E0, reflecting deployment without knowledge of the active error
structure. This separation prevents a structure-specific null from leaking
oracle information into the operational rule.

\subsection{Fixed-budget decision rules}\label{sec:rules}

R0 is the CI-only comparator. It certifies when the target gap interval excludes
zero. R1 allocates 10\% of the budget to the structure-designated probe and
certifies only when the target interval excludes zero and the probe does not
alarm. Let $W_f$ be the event that the reduced-shot target decision is wrong
and $A_f$ the alarm event. Define
\begin{equation}
  q_f=\Prb(W_f),\quad d_f=\Prb(A_f\cap W_f),\quad R_f=q_f-d_f.
\end{equation}
Against an all-target error probability $q_0$, the probe lowers wrong-decision
risk if and only if
\begin{equation}
  d_f>q_f-q_0.
  \label{eq:worth}
\end{equation}
The left side is useful interception; the right side is the error cost of
removing target shots. Equation~\eqref{eq:worth} is exact for the abstaining
policy and requires no independence assumption. Because abstention can
mechanically reduce recorded errors, the study reports false certification and
coverage separately rather than selecting a post hoc scalar loss.

\section{Experiment Setup}

\subsection{Graphs, candidates, and error structures}

The confirmatory holdout contains 100 independently generated graphs, four in
each width--depth stratum for widths 4--8 and depths 1--5. Graph families are
34 Erd\H{o}s--R\'enyi, 33 Watts--Strogatz, and 33 Barab\'asi--Albert. Each graph
has two frozen QAOA candidates and 32 scoring repeats per structure, probe, and
parameter condition. Candidates are compiled to \code{rz,sx,x,cx} with identity
layout, no routing, and optimization level zero. One candidate receives a
logically cancelling CX--CX pair after each cost layer, creating a frozen gate
count contrast without changing its ideal unitary.

Table~\ref{tab:errors} lists the seven non-null structures. E5 and E5$'$ retain
their independent marginal flips in the matched null; E6 retains its baseline
assignment rates. These matched nulls repair a v3 defect in which probes could
alarm on the baseline construction rather than on the tested correlation or
drift component.

\begin{table*}[t]
\centering
\caption{Each non-null structure is paired with a matched null that keeps its
baseline features and sets only the tested component to zero, so a probe
cannot score visibility by reacting to the baseline. E5$'$ is a synthetic
stress test, not a hardware-derived model.}
\label{tab:errors}
\begin{tabularx}{\textwidth}{@{}lXXl@{}}
\toprule
Code & Physical interpretation & Frozen strength grid & Matched null \\
\midrule
E1 & Two-qubit depolarization after every hard CX & $p\in\{0.0025,0.005,0.01,0.02,0.04\}$ & $p=0$ \\
E2 & Every tagged cost rotation runs as $R_Z((1+\epsilon)\theta)$ & $\epsilon\in\{\pm0.005,\pm0.01,\pm0.02,\pm0.04\}$ & $\epsilon=0$ \\
E3 & Coherent $R_{ZZ}(\zeta)$ immediately after every hard CX & $\zeta\in\{\pm0.005,\pm0.01,\pm0.02,\pm0.04\}$ & $\zeta=0$ \\
E4 & Independent asymmetric assignment error on every qubit & $(\Pr(1|0),\Pr(0|1))$: $(0.01,0.02)$, $(0.02,0.04)$, $(0.04,0.08)$, $(0.08,0.12)$, and the four swapped pairs & $(0,0)$ \\
E5 & Correlated pair flips on the first edge, marginal rate 0.10 & $\tau\in\{\pm0.0025,\pm0.005,\pm0.0075,\pm0.01\}$ & $\tau=0$, marginal flips kept \\
E5$'$ & Synthetic three-body flips on the first connected triple, marginal rate 0.25 & $\tau\in\{\pm0.00390625,\ldots,\pm0.015625\}$ & $\tau=0$, marginal flips kept \\
E6 & Eight-block assignment drift, $e_0=0.04+\delta t$, $e_1=0.06-\delta t$ & $\delta\in\{\pm0.005,\pm0.01,\pm0.02,\pm0.03\}$ & $\delta=0$, baseline kept \\
\bottomrule
\end{tabularx}
\end{table*}

E0 is a clean global null with equal deterministic offsets added to both
candidates. It cannot change a decision contrast. E7, fake-backend snapshots,
probe combinations, and additional diagnostics are exploratory and outside the
42-cell primary matrix.

\subsection{Preregistered matrix and endpoints}

The primary matrix comprises six probes by seven structures: 42 cells with 4
visible, 13 partial, and 25 blind predictions. Analytic classification yields
23 invariant and 19 sensitive cells (Table~\ref{tab:matrix}). The strata are
never pooled into one headline agreement score. All 23 invariant cells are
predicted blind; the sensitive stratum contains every visible and partial
prediction and two blind cells, P1/E6 and P3/E6, whose small block-dependent
population contrast is predicted to stay below the finite-budget power
boundary.

\begin{table}[t]
\centering
\caption{The preregistered matrix predicts blindness for most cells, and every
non-blind prediction lies in the budget-sensitive stratum. These are
predictions, not observations.}
\label{tab:matrix}
\begin{tabular}{@{}lrrrr@{}}
\toprule
Stratum & Visible & Partial & Blind & Total \\
\midrule
Invariant & 0 & 0 & 23 & 23 \\
Sensitive & 4 & 13 & 2 & 19 \\
\midrule
Total & 4 & 13 & 25 & 42 \\
\bottomrule
\end{tabular}
\end{table}

The primary estimand is
\begin{equation}
 \Delta_{\rm FC}=\frac{1}{7}\sum_{s=1}^{7}
 \{\mathrm{FC}_s(R0)-\mathrm{FC}_s(R1)\},
 \label{eq:dfc}
\end{equation}
where false certification means certifying the wrong ideal winner. A 95\%
interval is obtained from 10,000 stratified graph-cluster bootstrap resamples
with shared graph indices across rules and structures. The prespecified outcome
is ``improvement'' when the lower endpoint is positive, ``no decision problem''
when the 97.5th bootstrap percentile of pooled CI-only false certification is
below 0.01, and ``inconclusive'' otherwise.

Secondary endpoints include coverage loss, interception among CI-only false
certifications, behavior on flip and non-flip subsets, a structure-agnostic P2
rule, and visibility agreement within each analytic stratum. Readout
certificates are restricted to cells whose probe measurements address the
modeled readout component; they are not interpreted as certificates against
gate or state-preparation noise.

\subsection{Reproducibility and evidence separation}

All stochastic artifacts derive from a single NumPy \code{SeedSequence} key
containing split, artifact, structure, probe, width, depth, graph, candidate,
repeat, and parameter indices. Split 10 is development; split 20 is holdout;
split 99 is reserved for software tests and analytic design checks and is never
reported as physical or confirmatory evidence. Raw and failed rows are
append-only. Holdout access requires a second lock containing thresholds,
operating points, manifest hashes, and platform policy.

Protocol v4 and its complete execution code were registered on OSF at
\url{https://osf.io/r63vn/} at 2026-11-13 01:45:49 UTC, commit
\code{e3908dbdc9e73b3f5269195bf1839d74756e32e8}. It supersedes but does not
erase the public v3 record (OSF \code{j67wy}, commit \code{e8b5300}). This
manuscript snapshot uses repository HEAD \code{bdfeb5a}. At the snapshot date,
the Task 2-2 development computation had been executed but its outputs were not
yet validated or committed; none of them is opened, summarized, or used here.

\section{Results}

\subsection{A0: analytic and noiseless design results}

All results in this section are pre-data results: analytic derivations or
split-99 software-design checks. They are not development, holdout, hardware,
or physical-probe performance evidence.

QAOA--MaxCut has global bit-complement symmetry under the frozen ideal model,
so $\langle Z_i\rangle=0$. For uniform independent assignment noise with
$\widetilde Z_i=\lambda Z_i+\delta$, each edge contrast scales by $\lambda^2$;
hence
\begin{equation}
  g_{\rm noisy}=\lambda^2 g,\qquad
  \frac{\beta}{\Delta}=-(1-\lambda^2)\in(-1,0].
  \label{eq:readoutscale}
\end{equation}
Uniform E4-type readout noise therefore shrinks the gap but cannot reverse its
sign in the population limit. The same bound applies to the frozen E6 mixture.
This result limits what decision-flip evidence these structures can supply; it
does not establish that their probes are ineffective for detecting readout
change.

A noiseless density-matrix split-99 check over 50 graphs found no exact ideal
ties. At the strongest selected signs, median $\beta/\Delta$ values were
$-0.7358$ for E1, $-0.0912$ for E2, and $-0.0781$ for E3. The largest observed
fraction with $\beta/\Delta\le-1$ over the frozen gate grids was 0.12 (E2).
These are deterministic design diagnostics used to select and stress the
protocol; they are not estimates of holdout prevalence.

\subsection{A1: the registered v3 defect was detectable before noisy data}

After v3 registration and before any split-10 or split-20 row existed, an audit
proved five definition-level conflicts. P4a/E5, P4a/E5$'$, P4b/E5$'$,
P4a/E6, and P4b/E6 were predicted blind or partial, but their frozen statistics
necessarily responded to baseline marginal or assignment error when calibrated
against a perfectly clean E0. The corresponding visibility label would be
visible by construction. Eight P0 cells were also identified as high risk,
because target non-resolution is not a bias-visibility statistic; those eight
were projections, not analytic proofs.

The audit motivated three surgical changes in v4: P0 became a comparator rather
than a probe, every structure acquired a matched null, and each P4 cell received
one frozen component rather than an omnibus maximum. The original v3 files and
registration remain unchanged. This chronology is evidence of protocol repair,
not evidence in favor of the v4 scientific predictions.

\subsection{A2: semantic checks corrected four derivations}

Before v4 registration, exact expected-count calculations evaluated every cell
on 10 split-99 graphs of widths 4--6. The first pass disagreed with the analytic
invariant/sensitive stratum in four cells: P3/E2, P5/E3, P1/E6, and P3/E6.
Their derivations were corrected and both passes retained. A second pass over
the same design had zero disagreements across 42 cells. Because it followed
the corrections, this second pass is an internal consistency check, not an
independent validation of the predictions.

\subsection{A3: design power is limited by rare decision failures}\label{sec:power}

Power analysis v5 used only split-99 exact-population inputs and a frozen grid
of synthetic detection and false-alarm rates. It did not use visibility labels,
development outputs, holdout outputs, device parameters, or hardware data. In
the 50-graph design population, the strongest selected gate settings reversed
4 graphs for E1, 5 for E2, and 3 for E3; E4--E6 and E5$'$ reversed none. The
CI-only false-certification mean was 0.0200 for E2 and zero for the other six
structures at the modeled budget.

Across the preregistered design scenarios, the probability that the primary
analysis would declare ``no decision problem'' was approximately 0.685--0.775,
and the probability of an ``improvement'' conclusion did not exceed 0.15. The
calculation therefore predicts that the main limitation may be the scarcity of
false certifications to intercept, not merely probe sensitivity. This is a
prospective design warning, not a study result.

\begin{table}[t]
\centering
\caption{The manuscript intentionally stops at pre-data evidence. Development,
holdout, and hardware claims remain unavailable at this stage.}
\label{tab:status}
\begin{tabularx}{\columnwidth}{@{}p{0.29\columnwidth}X@{}}
\toprule
Evidence class & Status at snapshot \\
\midrule
Analytic theory & Available; assumptions stated \\
Split-99 checks & Available; design/software evidence only \\
OSF protocol v4 & Registered with execution code \\
Split-10 development & Executed; not yet validated or committed; excluded \\
Split-20 holdout & Not accessed \\
Hardware data & None \\
\bottomrule
\end{tabularx}
\end{table}

\section{Conclusion}

Protocol v4 turns a broad question---whether a runtime diagnostic ``sees''
noise---into a fixed-budget decision test. It separates analytic blindness from
finite-shot failure, uses matched nulls, and evaluates whether alarms intercept
wrong certifications rather than merely reacting to an error channel. Before
any reported development or holdout result, the project has established an
auditable 42-cell prediction matrix, repaired a registered protocol defect,
and shown prospectively that rare decision flips may dominate statistical
power. It has not established that any probe improves decisions, generalizes to
hardware, or detects unknown noise. Those conclusions remain contingent on the
locked development and confirmatory stages.

\section{Discussion}

\subsection{What the present evidence supports}

The analytic results support three bounded claims. First, sampling confidence
and differential-bias control are distinct. Second, probe visibility is
mechanism-specific: a statistic can be exactly invariant to a harmful direction
even when another probe responds. Third, a responsive probe is not necessarily
worth its shots; decision value requires Eq.~\eqref{eq:worth}, not high alarm
power alone. The documented v3 failure also shows why these claims require
matched nulls and a strict separation between comparator uncertainty and bias
alarms.

\subsection{Hypotheses and possible explanations}

The following are expectations to be tested, not findings. If the primary
endpoint returns ``no decision problem'', the most direct explanation is the
design calculation of Section~\ref{sec:power}: at this budget, CI-only false certification
is concentrated in a few graphs under one structure, leaving little for any
probe to intercept. If sensitive cells predicted partial are observed blind,
a plausible explanation is that the population contrast is real but smaller
than the frozen power boundary allows at $n_P=1{,}638$. Disagreements in the
invariant stratum would instead indicate an implementation or derivation
error, because their population statistics are equal by construction.

\subsection{Limitations}

The error classes are structured simulations, not a representative sample of
real devices. E5$'$ is synthetic and has no claimed hardware evidence. E4 and
the frozen E6 cannot reverse the population decision under the idealized
symmetry assumptions of Eq.~\eqref{eq:readoutscale}; they test visibility and
calibration, not a rich source of false certifications. Candidate generation,
graph families, QAOA depths, and the deliberate CX--CX contrast limit external
generalization. Partial-identification language does not remove assumptions
about the admissible mechanism class. Moreover, semantic tests performed on the
same design after corrections are not independent replication. The matched
nulls are defined inside the simulated error model, so visibility labels
depend on that model. The individual lemmas rest largely on established
statistics and information theory; the contribution is their joint,
preregistered use, not new bounds. Hardware validation has not been
authorized under the current protocol.

The primary endpoint may be under-informative if false certifications are rare.
In that case, a ``no decision problem'' outcome is scientifically different
from evidence that probes work. Conversely, an inconclusive result cannot be
converted post hoc into a positive visibility claim. Coverage and abstention
cost must accompany every reduction in false certification.

\subsection{Reproducibility and next evidence gates}

The confirmatory lock must be written before split-20 access and must freeze the
development-selected operating points, thresholds, hashes, and execution
platform. Public continuous integration should reproduce the test suite and
the deterministic reference phase on clean Windows and Linux runners, recording
package and BLAS configurations. Platform differences should be retained and
reported rather than hidden behind relaxed tolerances. Only after these gates
and the holdout analysis should a hardware micro-study be considered. Such a
study would be an external anchor, not retroactive validation of the simulated
error family.

\bibliographystyle{abbrv}
\bibliography{references}

\appendix

\section{Blindness lemma}

Let $M$ denote an admissible noise mechanism, $Q_M$ the joint population law
of all target and probe observations, and $d(M)\in\{-1,+1\}$ the ideal decision.
If two admissible mechanisms $M_+$ and $M_-$ satisfy
\begin{equation}
  Q_{M_+}=Q_{M_-},\qquad d(M_+)\neq d(M_-),
\end{equation}
then no rule measurable with respect to the recorded observations can be
uniformly correct over the model class. Indeed, any rule has the same output
distribution under the identical laws, so its error probabilities under the
opposite decisions sum to one. At least one is therefore at least $1/2$. This
is an identification statement, not a finite-sample approximation.

\section{Finite-shot necessary condition}

For two data-generating states $P_0$ and $P_1$ that require opposite binary
decisions, any test based on $N$ independent observations obeys
\begin{equation}
  P_0(\text{error})+P_1(\text{error})
  \ge 1-\operatorname{TV}(P_0^{N},P_1^{N}).
\end{equation}
Using $\operatorname{TV}(P,Q)\le H(P,Q)$ and
$H^2(P_0^N,P_1^N)\le N H^2(P_0,P_1)$ under the convention
$H^2(P,Q)=\int(\sqrt p-\sqrt q)^2$, a necessary condition for both errors to
be at most $\alpha<1/2$ is
\begin{equation}
  N\,H^2(P_0,P_1)\ge (1-2\alpha)^2.
\end{equation}
The bound is pairwise and necessary; it neither constructs an optimal probe nor
replaces the exact protocol calibration.

\section{Local information is not global sign identification}

Two analytic counterexamples separate a local Fisher null-space condition from
global sign identification. For $X\mid\theta\sim\mathcal N(\theta^2,1)$ and
$g(\theta)=\theta$, Fisher information is positive at every $\theta\neq0$, yet
$P_a=P_{-a}$ and the two equivalent points have opposite decisions. Conversely,
for $X\mid(u,v)\sim\mathcal N(u,1)$ with $v\in(-1/2,1/2)$ and
$g_m(u,v)=m(1+v)$, the unobserved $v$ direction changes the gap magnitude and
violates the local null-space condition, but the full identification region
$(m/2,3m/2)$ remains positive. Local sensitivity and global decision blindness
therefore do not imply one another without additional regularity and
localization assumptions.

\section{Fixed-budget Proposition 2}

Fix a total budget $N$, give a fraction $f\in[0,1)$ to the probe and
$(1-f)N$ to the target, and return \emph{unresolved} whenever the probe alarms.
With $W_f$, $A_f$, $q_f$, and $d_f$ as in Section~\ref{sec:rules}, the wrong-decision
probability is exactly
\begin{equation}
  R_f=\Prb(W_f\cap A_f^{c})=q_f-d_f,
\end{equation}
and the all-target baseline is $R_0=q_0$. Hence $R_f<R_0$ if and only if
$d_f>q_f-q_0$, which is Eq.~\eqref{eq:worth}. Writing
$\eta_f=\Prb(A_f\mid W_f)$ for $q_f>0$ gives the equivalent condition
\begin{equation}
  \eta_f>1-q_0/q_f.
\end{equation}
No Gaussian or independence assumption is used. Under an additional Gaussian
model with independent probe and target data, $q_f=\Phi(-m\sqrt{(1-f)N}/\sigma)$
for noisy gap $m=\Delta+\beta>0$, and $\eta_f$ equals the probe's alarm
probability; this approximation does not apply to P5, P1b, or GVS, which reuse
target counts. If unresolved outcomes carry loss $\lambda$, the probe lowers
$R_f+\lambda\Prb(A_f)$ if and only if $d_f>q_f-q_0+\lambda\Prb(A_f)$. The study
reports false certification and coverage separately instead of fixing
$\lambda$ post hoc. When $m<0$, removing target shots moves $q_f$ toward $1/2$
and can lower $q_f$ mechanically; that reduction is not evidence of
visibility.

\section{Protocol correction chronology}

\begin{table}[h]
\centering
\caption{Protocol changes were made before split-10 or split-20 evidence was
used. Earlier registrations remain public rather than being overwritten.}
\begin{tabularx}{\columnwidth}{@{}p{0.20\columnwidth}X@{}}
\toprule
Stage & Public record and scientific role \\
\midrule
v2 review & Pre-development review proved that uniform readout shrinks but does not reverse the frozen QAOA decision; it motivated structure-specific operating points. \\
v3 & OSF \code{j67wy}, commit \code{e8b5300}; retained unchanged after a pre-data audit found P0-role and null-calibration defects. \\
v4 & OSF \code{r63vn}, commit \code{e3908db}; registered matched nulls, 42-cell strata, the decision endpoint, and complete execution code. \\
Snapshot & Repository \code{bdfeb5a}; Task 2-2 uncommitted development outputs are excluded from this manuscript. \\
\bottomrule
\end{tabularx}
\end{table}

\end{document}
